\section{Ablation}\label{sec:ablation}

The calibrated predictor has four components that the study can switch one at a time: the scalars (given at the truth,
given at their shipped values, or learned), the discrepancy prior (Kennedy and O'Hagan, orthogonal or scaled), the
discrepancy itself (kept or dropped), and the form of the engine (affine or curved capacity). Table~\ref{tab:ablation} gives
the RMSE of each variant on six cells; the capacity form is the subject of Figure~\ref{fig:capacity} and is not repeated.

\begin{table}[!t]
\caption{Ablation of the calibrated predictor: RMSE (kW for power, t/h for throughput) on six cells with the scalars at the
truth, at their shipped values and learned under the three discrepancy priors, with the discrepancy dropped (the engine refit
alone), and least squares on the setting and the feed size for reference; \repsPerCell{} replicates per cell.}
\label{tab:ablation}
\centering\footnotesize
\resizebox{\columnwidth}{!}{\begin{tabular}{@{}lrrrrrrr@{}}
\toprule
Cell & True & Shipped & KO & Orth. & Scaled & No discr. & LS \\
\midrule
V0 power, surveys, 10 & 45.2 & 46.9 & 46.1 & 71.4 & 46.1 & 64.3 & 46.9 \\
V0 power, surveys, 160 & 39.2 & 39.5 & 39.4 & 49.3 & 39.4 & 55.4 & 39.5 \\
V0 tph, surveys, 10 & 168.6 & 168.6 & 169.5 & 169.4 & 167.4 & 183.5 & 174.4 \\
V2 tph, surveys, 10 & 168.1 & 168.2 & 169.0 & 168.9 & 167.1 & 186.2 & 173.7 \\
V2 power, extrap., 40 & 45.4 & 46.1 & 45.8 & 53.5 & 45.5 & 50.7 & 45.9 \\
V1($-1$) power, wide, 40 & 41.8 & 42.0 & 41.6 & 93.4 & 41.9 & 94.9 & 40.7 \\
\bottomrule
\end{tabular}
}
\end{table}

\subsection{The scalars}

Learning the scalars costs little against knowing them. On the power of V0 at ten surveys the predictor with the true scalars
reaches \rVzeroSTenKwOr~kW, the shipped scalars \rVzeroSTenKwNom{} and the learned ones \rVzeroSTenKwKo; at 160 surveys the
three are within half a kilowatt of one another and of least squares. On throughput the scalars hardly matter at all, because
the engine's throughput is affine and its scale is absorbed by the plane: the true and the shipped scalars give the same
\rVzeroSTenTphOr~t/h. The cost of learning the scalars is largest at ten surveys on the real design (\rVzeroRealKwKo{} against
\rVzeroRealKwOr~kW), where ten observations per output are shared between eight scalars, three discrepancy coefficients
and the noise; it is the price of the uncertainty that the band of the method carries and the residual band of a refit does not.
The learned predictor here is the linearized one, which fixes the noise scale at the plant's and drops the truncation of the
priors; Section~\ref{sec:results-sampler} measures the full posterior on the same replicates.

\subsection{The discrepancy prior}

The orthogonal prior is the costly change. With the product's proper priors on the scalars it raises the power RMSE on V0 from
\rVzeroSTenKwKo{} to \rVzeroSTenKwOrth~kW at ten surveys and from \rVzeroSBigKwKo{} to \rVzeroSBigKwOrth{} at 160, and on V1($-1$)
over the wide window to \rVamWFortyKwOrth~kW, next to the engine without a discrepancy (\rVamWFortyKwRefit). The mechanism is
Proposition~\ref{prop:priors}(a): the eight sensitivities of power on the training surveys span most of the directions of the
three discrepancy regressors, so the orthogonal complement leaves the discrepancy little room, and the plane the plant needs
must be carried by the scalars, whose priors forbid it. In the diffuse limit the two priors predict alike, by part (b); with the
priors the product uses, they do not, and the orthogonal prior behaves like the engine without a discrepancy. On throughput,
where only the capacity base has a sensitivity, the two priors are within a tenth of a percent of each other. The scaled prior
of Gu and Wang changes little on power on the surveys window (it is \scaledWideWorse\pct{} worse than the Kennedy and O'Hagan prior on
the wide window at ten surveys) and is the best of the three on throughput at ten surveys (\rVzeroSTenTphSc~t/h against
\rVzeroSTenTphKo{} for the Kennedy and O'Hagan prior), the shrinkage it adds along the well-measured directions being where the
throughput's plane mismatch is small. H4 as stated, that the priors change the risk less than they change the meaning of the
scalars, failed for the orthogonal prior: at proper priors it changes the risk most.

\subsection{The discrepancy}

Dropping the discrepancy is the largest change of all, in both directions. On every plant whose form the engine does not share
the refit engine is the worst predictor of the table except the orthogonal prior on two power rows, up to \rVamWFortyKwRefit~kW on V1($-1$) against \rVamWFortyKwLs{} for
least squares, and its band covers as little as \covRefitMin{} of the test points. On the capacity study with the engine of the
plant's own form it is the best predictor of the study (Section~\ref{sec:results-capacity}). The discrepancy plane therefore
buys robustness against a wrong form at the price of the variance of three coefficients, and Theorem~\ref{thm:fixed} prices
it: at most $3\sigma^2/n$ when the form is right, and the whole curvature of the plant when the form is wrong. The second price
is the larger on every plant of the study whose form the engine does not share, and the first fades with the surveys, so on
the evidence of this study the discrepancy should be kept unless the engine's form has been validated over the operating range
by other means; ten surveys of throughput and power cannot validate it.

\subsection{What the ablation says about the method}

Read together, the variants say that CrusherBayes is close to the best predictor available to it in every cell of the study
except one class: when the engine's form is exactly right, the discrepancy should be dropped, which no set of surveys can
establish in advance. Its band covers at about the nominal rate, as the Student-$t$ band of least squares does, and is the
narrower of the two in every cell of Table~\ref{tab:main}; its accuracy is that of least squares within the curvature and plane
terms of Theorem~\ref{thm:fixed}; and with learned scalars its RMSE is \rVzeroSTenKwKo~kW on ten random surveys against
\rVzeroSTenKwOr~kW with the true scalars, and \rVzeroRealKwKo{} against \rVzeroRealKwOr~kW on the real design. The choice of prior matters more than the choice of discrepancy scale, and the prior that the calibration literature
recommends for the meaning of the scalars is the wrong one for prediction with this engine.
